\documentclass[11pt]{article}
\usepackage[final]{acl}
\usepackage{times}
\usepackage{latexsym}
\usepackage[T1]{fontenc}
\usepackage[utf8]{inputenc}
\usepackage{microtype}
\usepackage{inconsolata}
\usepackage{graphicx}
\usepackage{booktabs}

\title{Is That for Me? Addressee-Aware Turn-Keeping for Open Full-Duplex Speech Models under Same-Voice Side-Talk}

\author{Zimo Fan \and Ruihuan Gao \and Ruize Sun \and Weiyi Sun \\
  ECE 6256 Conversational AI (Fall 2026), School of ECE, Georgia Institute of Technology \\
  \texttt{\{zfan321,rgao302,rsun98,wsun377\}@gatech.edu} \\
  \textbf{Anchored in:} Unit~6 Voice (L18--L21) and Unit~7 Multi-agent--human (L22--L23)}

\begin{document}
\maketitle

\begin{abstract}
Open full-duplex speech models treat every voice in the user channel as meant for them, so on three recent benchmarks they answer when the user turns to a third person. The shipped fix in Gemini 3.8 Live is unpublished, and open remedies gate on speaker identity, which fails when the side-talk is in the user's own voice. We measure how much identity gating solves, build a de-shortcut same-voice test set with voice cloning, LoRA-fine-tune Moshi on cue-controlled synthetic side-talk, and ablate lexical against acoustic cues.
\end{abstract}

\section{Problem and gap}
A full-duplex model must decide, frame by frame, whether incoming speech is for it. A 2026 survey of full-duplex systems \citep{lu2026survey} lists ``third-party / not-directed speech'' as an ontology cell that speaker-embedding and relevance gates target but have ``not shown solved''. On Full-Duplex-Bench (FDB) v1.5 \citep{lin2025fdb15}, when the user talks to someone else, GPT-4o Realtime \citep{openai2024gpt4o} responds on 91\% of clips and Freeze-Omni \citep{wang2024freezeomni} on 58\%; Moshi responds on 20\% and is unclassifiable on another 59\%, and the authors conclude that systems ``must strengthen addressee detection''. The HumDial challenge \citep{wang2026humdial} and MSI-Bench \citep{xiong2026msibench} report the same failure: ``models often respond even when no one addressed them''. Gemini 3.8 Live ships always-on ``proactive audio'' that ``can proactively decide not to respond'' \citep{google2026gemini38live}, and Realtime-Venus \citep{antgroup2026venus} reports 88\% correct ignoring on this cell; neither releases weights or a method.

Every open remedy gates on \emph{who} is speaking: FireRedChat's personalised VAD \citep{chen2025fireredchat,ding2019pvad} and IRAF's ECAPA-embedding gate inside an end-to-end model \citep{zhong2026iraf}. Yet FDB v1.5 deliberately separates two cells: \emph{bkg\_speech} is a different voice at $-15$~dB, whereas \emph{talking\_other} is the \emph{user's own voice} at $-8$~dB with far-field reflections, saying something meant for a third person. Identity gating can only address the first. No public model has been trained on same-voice side-talk and no corpus labels it at scale (HumDial's training set has zero such items).

\section{Research questions and hypotheses}
\textbf{RQ1} Which part of the not-addressed cell does identity gating solve? \textit{H1:} a verification gate mutes most different-voice speech but passes most same-voice side-talk. \textbf{RQ2} Can a Moshi-class model learn to keep its turn under same-voice side-talk while still yielding to genuine interruptions? \textit{H2:} LoRA on cue-controlled synthetic side-talk at least halves base Moshi's \emph{side-talk onset rate} (share of side-talk events after which the model starts speaking, by VAD) at $\leq 5$ points more missed interruptions. \textbf{RQ3} Which cue does it learn? \textit{H3:} with acoustic cues (level drop, far-field colouring) present, the gain survives removing the lexical cue but not the acoustic one.

\section{Proposed system}
\textbf{Base and baselines.} Moshiko-7B \citep{defossez2024moshi} (CC-BY~4.0, official \texttt{moshi-finetune} LoRA) is the fine-tuning target; PersonaPlex-7B \citep{roy2026personaplex} and Freeze-Omni are inference baselines. Two gates over the same Moshi aim at side-talk directly: a vocative rule (the pilot's comma-within-three-words rule) and a cascade addressee classifier (streaming ASR plus a local 8B LLM judging whether each utterance addresses the system).
\textbf{G0 identity gate (no training).} ECAPA \citep{desplanques2020ecapa} cosine between each 0.5--1~s user-channel window and an enrollment from the user's earlier speech; windows below a held-out-tuned threshold are zeroed before the model hears them.
\textbf{Data pipeline.} Behavior-SD \citep{lee2025behaviorsd} (108k synthetic dialogues, CC-BY~4.0) supplies base dialogues. An LLM writes side-talk addressed to a named third party (with and without a vocative) and genuine interruptions; Chatterbox \citep{resemble2025chatterbox} voice cloning renders them in the user's voice; they enter the user channel under three acoustic conditions: FDB-matched far-field ($-8$~dB, high-shelf cut, reflections), near-field (lexical cue only) and content-neutral far-field (acoustic cue only). Different-voice background speech (LibriSpeech \citep{panayotov2015librispeech}, 0--15~dB SNR) is the control cell.
\textbf{Training target.} Moshi keeps producing its utterance through non-addressed speech and yields on addressed interruptions; LoRA rank 128, about 2{,}000 steps, one H100 ($\approx$40~GB).

\section{Experimental design}
\textbf{Test sets.} (i) FDB v1.5: talking\_other, bkg\_speech, backchannel, interruption (498 clips). (ii) HumDial-FDBench test set (600 third-party and 200 talking-to-others items, human actors). (iii) Our de-shortcut set: same vs.\ different voice $\times$ acoustic cue $\times$ vocative, on unseen voices, with far-field genuine interruptions and ``Hey Moshi, \ldots'' vocatives aimed at the system. (iv) A small human-recorded set (50--100 clips). (v) Multi-turn sessions: 200 held-out Behavior-SD dialogues of 1--2~min, each with 2--4 same-voice side-talk events and 1--2 genuine interruptions.
\textbf{Metrics} (conditions: Moshi, +vocative rule, +ASR-LLM gate, +G0, +LoRA, +G0+LoRA, PersonaPlex, Freeze-Omni)\textbf{.} Side-talk onset rate (primary), with FDB respond labels alongside; response rate and stop latency on genuine interruptions; content quality (FDB judge); gate latency within the 80~ms frame. On set (v), per-session scores: share of sessions with every event handled correctly, cumulative spurious responses, and resumption (LLM-judged: does the model return to its prior topic after side-talk).
\textbf{Estimands and analysis.} Primary: relative change in side-talk onset rate (LoRA vs.\ base) with the interruption miss rate held within 5 points. Secondary: G0's split between different-voice and same-voice speech. Paired bootstrap CIs by clip (by session on set v); three LoRA seeds; the $2\times2$ cue ablation on set (iii) answers RQ3.

\section{Pilot evidence}
On FDB v1.5 (498 clips, ECAPA) the identity cue behaves as H1 predicts (Table~\ref{tab:pilot}): the gate separates background speech from real interruptions almost perfectly (AUC 0.9998) yet accepts 88--92\% of same-voice side-talk at thresholds that keep 95\% of the user's own speech. The pilot also exposed two shortcuts: 99 of 100 side-talk clips open with a vocative, so a comma-within-three-words rule flags 99\% at 3\% false positives, and all side-talk is far-field while all interruptions are near-field. Caveat: talking\_other, bkg\_speech and backchannel share one TTS user voice while interruptions use many; matched on voice, ECAPA separates the user's clean speech from side-talk at AUC only 0.61, another reason for set (iii).

\begin{table}[t]
\centering\small
\begin{tabular}{lccc}
\toprule
FDB v1.5 scenario & cos.\ & pass & muted \\
\midrule
user\_interruption (keep) & .49 & .995 & .13 \\
talking\_other (want mute) & .39 & \textbf{.89} & \textbf{.27} \\
user\_backchannel (keep) & .19 & .52 & .47 \\
bkg\_speech (mute) & .06 & .01 & .83 \\
\bottomrule
\end{tabular}
\caption{Identity cue on FDB v1.5: median ECAPA cosine to the user; clips accepted as the user at the EER threshold; audio muted by the 0.5~s G0 gate.}
\label{tab:pilot}
\end{table}

\section{Plan, roles and feasibility}
Week~1 (to 11/12): baselines and side-talk gates on FDB v1.5 and HumDial, G0, LoRA smoke test, synthesis v0. Milestone 11/19: baseline table with the G0 split, first LoRA. Weeks 3--4: LoRA sweep over cue conditions and data size, de-shortcut, human and multi-turn sets, latency. One member each on evaluation and baselines, synthesis and test sets, G0 and analysis, LoRA. \textbf{Compute (PACE-ICE, estimated).} Inference on an L40S ($\approx$24~GB): $\sim$15~h of audio per condition $\times$ $\sim$12 runs $\approx$ 150 GPU-h. LoRA on one H100: $\approx$3~h $\times$ $\sim$20 runs $\approx$ 60 GPU-h. API (GPT-4o FDB judge, resumption judge, side-talk generation): about \$150.

\section{Risks and fallback}
(1) The model may learn ``quieter means ignore'' rather than addressee semantics; the cue ablation turns this into a finding. (2) \texttt{moshi-finetune} is validated only on Moshiko; a week-1 smoke test gates the rest. (3) Synthetic side-talk may not transfer to real rooms; HumDial's human recordings and our recorded set bound the gap. If LoRA fails, the deliverable is a measurement paper: every non-LoRA condition on all five test sets, the identity-gate split and the de-shortcut analysis.


\bibliography{proposal}

\end{document}
